\section{Fonaments de planificació de tasques}

La planificació decideix quina tasca pendent rep el recurs disponible. En el
model del TFM la decisió és dinàmica, perquè poden arribar tasques mentre altres
s'executen, i en línia, perquè només es coneix l'estat disponible en cada
instant. La política no crea capacitat: modifica com es distribueix l'espera
entre tasques i classes.

\subsection{FIFO}

FIFO tracta l'antiguitat d'arribada com a únic criteri. És fàcil d'explicar,
implementa un desempat estable i serveix com a línia base. El seu límit apareix
quan una tasca urgent arriba darrere d'un conjunt ampli de treballs o quan una
tasca llarga ocupa el primer lloc. Tot i això, no s'ha de descriure FIFO com una
política formalment justa sense indicar quina definició d'equitat s'utilitza.

\subsection{Prioritat estricta}

La prioritat estricta selecciona primer la tasca amb la prioritat més alta i
utilitza habitualment FIFO per desempatar. Permet donar una resposta preferent
a una classe urgent, però una arribada sostinguda de tasques d'alta prioritat
pot ajornar indefinidament les de nivell inferior. RabbitMQ tracta aquest risc
d'inanició i mostra que el comportament concret depèn del tipus de cua i de la
implementació \parencite{rabbitmqPriority}.

La càrrega és decisiva. Wierman i Harchol-Balter mostren que el comportament
d'equitat d'una disciplina pot variar segons la càrrega i la distribució dels
temps de servei \parencite{wierman2003unfairness}. Wierman també assenyala que
l'estudi tradicional del temps de resposta, la longitud de cua i el throughput
s'ha de complementar amb alguna noció d'equitat \parencite{wierman2011fairness}.
Encara que els models no coincideixen exactament amb el TFM, justifiquen provar
diversos nivells de saturació i observar els resultats per classes.

\subsection{Prioritat amb envelliment}

L'envelliment modifica la prioritat efectiva en funció del temps acumulat a la
cua. Una tasca de prioritat baixa pot ascendir gradualment fins a competir amb
tasques inicialment més urgents. El mecanisme necessita concretar una fórmula,
un interval o un llindar de promoció; aquesta elecció forma part del model.

Becker i Schüle avaluen un mecanisme adaptatiu d'envelliment combinat amb sis
algorismes dinàmics. Els seus experiments redueixen l'espera de treballs de
prioritat baixa i milloren el makespan mitjà en part dels casos, però també
augmenten l'espera de treballs prioritaris i cap algorisme domina totes les
proves \parencite{becker2020aging}. El domini és una arquitectura heterogènia,
de manera que els valors no es poden extrapolar. La conclusió transferible és
que l'envelliment introdueix un compromís que cal mesurar.

\begin{table}[htbp]
  \centering
  \caption{Síntesi de les polítiques previstes}
  \label{tab:politiques}
  \begin{tabularx}{\textwidth}{@{}p{2.4cm}X X X@{}}
    \toprule
    Política & Criteri & Avantatge a contrastar & Risc principal \\
    \midrule
    FIFO & Ordre d'arribada & Simplicitat i línia base estable & Bloqueig de cap de línia \\
    Prioritat estricta & Prioritat inicial; FIFO en empat & Resposta preferent per a tasques urgents & Inanició de prioritats baixes \\
    Prioritat amb envelliment & Prioritat inicial i espera acumulada & Reducció d'esperes extremes & Sensibilitat a la fórmula i penalització de classes altes \\
    \bottomrule
  \end{tabularx}
\end{table}
